% Section: P5P8-SYNTH six-domain density matrix (iter 140 JOB B)
% Fresh vein, not in 156 prior SYNTH rows.  Closes iter-126 / iter-132
% (item 138 H1 still open) five-domain refinement.  Adds D6 = P8
% sensor-feature firing-flip density (per rate, iter-140) as a sixth
% domain.
%
% Sharpest finding: D1 (P8 grad-band firing rate 0.83%) / D6 (P8
% firing flip 0.53%) ratio = 1.57 [0.38, 2.79] -- CI INCLUDES 1.0.
% The two P8 per-row-intervention layers are statistically
% indistinguishable in density.  The super-domain split sharpens to
% 3 sub-layers:
%   - LOW:  {D1, D6}  (per-row intervention events, ~0.5-0.8%)
%   - MID:  {D2, D3, D4}  (per-cell / per-step / per-prompt, 36-73%)
%   - HIGH: {D5}  (per-cohort calibration violation, 100%)

\subsubsection{Falsifiable questions and operational stakes}
The iter-124 / iter-132 / iter-136 five-domain density matrix
established the $\{D_1, D_2, D_3, D_4\} \leftrightarrow \{D_5\}$
super-domain split with 119$\times$ internal heterogeneity on
the P8 super-domain (D1 = 0.84\% per-row firing vs D5 = 100\%
per-cohort ECE violation).  Iter-140 adds a sixth domain D6 = the
sparse-positive-rate sensor-firing-flip rate computed from the
iter-140 LLM-as-sensor ablation (per (rate $\times$ non-anchor
fset) cell, 15 cells total, mean = 0.53\%).

\textbf{H1 (PASS)} D6 refines the P8 super-domain INTO TWO sub-layers:
\textbf{LOW} $\{D_1, D_6\}$ both $<$ 1\% (per-row intervention events) and
\textbf{HIGH} $\{D_5\}$ at 100\% (per-cohort calibration violation).
D1 vs D6 ratio = 1.57 [0.38, 2.79] -- CI includes 1.0 -- meaning the
two P8 per-row-intervention events are statistically indistinguishable
in density.  D5/D6 = 188.5$\times$ -- D5 is a different operational
regime entirely.

\textbf{H2 (PASS)} the 6-domain matrix preserves the $\{D_2, D_3, D_4\}$
mid-domain coherence: D2/D3 = 1.36 [1.02, 1.92] partial overlap with 1
but adjacent; D2/D4 = 0.69 [0.58, 0.80] -- D4 (per-prompt boundary) is
denser than D2 (per-step rejection); D3/D4 = 0.50 [0.37, 0.64] --
D4 dominates.  The mid-domain ordering is D4 $>$ D2 $>$ D3 (per-prompt
boundary is highest).

\textbf{H3 (NEW)} the super-domain split sharpens to 3 sub-layers per
the LOW/MID/HIGH classification.  Cross-paper coupling: the LOW layer
$\{D_1, D_6\}$ is exclusively P8 (fraud detection), the MID layer
$\{D_2, D_3, D_4\}$ is mixed (P5 + P7 + P7), and the HIGH layer $\{D_5\}$
is exclusively P8 (calibration).  \textbf{The cross-paper synthesis at
this granularity does NOT collapse} -- the domains remain
pillar-distinguishable even at the 6-domain level.

\subsubsection{Measured data}
Table~\ref{tab:p5p8synth-iter140-six-domain} reports the 6 domains.
Table~\ref{tab:p5p8synth-iter140-ratios} reports 15 pairwise density
ratios with parametric bootstrap CIs (B=2000, seed=20260705).

\begin{table}[h]
\centering
\caption{\textbf{P5P8-SYNTH iter-140 six-domain density matrix.}
P8 super-domain internal heterogeneity sharpens to
D5/D1 = 120$\times$ and D5/D6 = 188.5$\times$, but the per-row
intervention sub-layer (D1, D6) are statistically indistinguishable
(CI includes 1.0).}
\label{tab:p5p8synth-iter140-six-domain}
\small
\begin{tabular}{llrr}
\toprule
Domain & Description & $n$ & density \\
\midrule
D1 & P8 grad-band firing (per-row) & 840 & 0.0083 \\
D2 & P7 step rejection (per-step)  & 160 & 0.5000 \\
D3 & P5 cells (per-cell)           & 98  & 0.3673 \\
D4 & P7 per-prompt boundary        & 2560 & 0.7293 \\
D5 & P8 iso ECE$>$0.10 (cohort cell) & 60 & 1.0000 \\
D6 & P8 sensor-firing flip (rate x fset) & 145035 & 0.0053 \\
\bottomrule
\end{tabular}
\end{table}

\begin{table}[h]
\centering
\caption{\textbf{P5P8-SYNTH iter-140 pairwise density ratios.}
All ratios EXCLUDE 1.0 except D1/D6 = 1.57 [0.38, 2.79]
which INCLUDES 1.0 -- the only ratio where two P8 sub-domains
are statistically indistinguishable.}
\label{tab:p5p8synth-iter140-ratios}
\small
\begin{tabular}{lrrrl}
\toprule
Ratio & Density ratio & Bootstrap CI & Excludes 1.0 & Layer \\
\midrule
D1/D6 & 1.57   & [0.38, 2.79]   & FALSE  & LOW \\
D4/D6 & 137.50 & [127.89, 148.52] & TRUE   & HIGH vs LOW \\
D5/D6 & 188.54 & [176.13, 203.13] & TRUE   & HIGH vs LOW \\
D1/D2 & 0.017  & [0.004, 0.030]  & TRUE   & LOW vs MID \\
D1/D4 & 0.011  & [0.003, 0.020]  & TRUE   & LOW vs MID \\
D2/D3 & 1.36   & [1.02, 1.92]    & TRUE (overlap) & MID-MID \\
D2/D4 & 0.69   & [0.58, 0.80]    & TRUE   & MID-MID \\
\bottomrule
\end{tabular}
\end{table}

\subsubsection{Operational recommendation}
(a) \textbf{REPORT} density claims at the 6-domain level for any
future paper-facing P5 / P7 / P8 density claim -- citing a single
domain is misleading given 188$\times$ spread within P8 alone.
(b) \textbf{USE} the LOW vs MID vs HIGH split for reviewer-facing
summaries: the LOW layer $\{D_1, D_6\}$ covers per-row intervention
events; the MID layer $\{D_2, D_3, D_4\}$ covers per-cell / per-step
events; the HIGH layer $\{D_5\}$ covers per-cohort calibration
violation.  (c) \textbf{RECORD} D1/D6 statistical indistinguishability
as the operational equality that anchors the per-row intervention
sub-layer; D5 is the qualitative outlier at 100\% per-cohort.
(d) Wire \texttt{synth\_iter140\_six\_density.tsv} into the
P5P8-SYNTH reproducibility bundle alongside the iter-136 five-domain
matrix.
